\scanpage{91}


i.e. 

\[ c=c_1(b_2^2-4c_2)+c_2(b_1^2-4c_1)+4c_1c_2\qquad\text{i.e.} \]

\[ \boxed{c=c_1b_2^2+c_2b_1^2-4c_1c_2}\ . \]

Comme on a de plus

\[ \boxed{b=b_1b_2}, \]

cela suffit à déterminer $Q$, comme

\[ Q(\underset{\substack{\rotatebox{-90}{$\in$}\\\Gamma O}}{\lambda}x_1*x_2\deleted{\unsure{)}}+\underset{\substack{\rotatebox{-90}{$\in$}\\\Gamma(L_1\otimes L_2)}}{u})=u^2+\lambda bu+\lambda^2c \]

\uline{Corollaire} Si $\delta$ est le discriminant de $Q$, on a

\[ \delta=\delta_1\delta_2\qquad\text{(en plus de }b=b_1b_2\text{)} \]

En effet

\[ \begin{aligned}\delta=b^2-4c&=b_1^2 b_2^2-4(c_1b_2^2+c_2b_1^2-4c_1c_2)\\ &=(b_1^2-4c_1)(b_2^2-4c_2)\end{aligned} \]

\uline{NB} Les formules simples à retenir, qui déterminent complètement $Q$ en termes de $Q_1$ et $Q_2$, quand on a des scindages, sont

\[ \boxed{\begin{aligned} b&=b_1b_2\\\delta&=\delta_1\delta_2\end{aligned}} \]

on en conclut pratiquement la formule pour $c$, puisque

\[ \delta=b^2-4c=\delta_1\delta_2\qquad\text{d'où, comme }b=b_1b_2 \]

\[ \begin{aligned}-4c=\delta_1\delta_2-b^2&=(b_1^2-4c_1)(b_2^2-4c_2)-b_1^2\cdot b_2^2\\ &=\unsure{4|4_1}(-b_1^2c_2\unsure{-|+}c_1b_2^2+4c_1c_2)\end{aligned} \]
% archive sheet number: 90
